\providecommand{\BookProfile}{screen}
\def\BookPrintProfile{print}
\newif\ifbookprint
\ifx\BookProfile\BookPrintProfile\bookprinttrue\else\bookprintfalse\fi

\documentclass[fontset=none,openany,10pt,twoside]{ctexbook}
\usepackage{graphicx}       % 插入图片
\graphicspath{{figure/}}
\usepackage{subcaption}     % 并排子图
\PassOptionsToPackage{hyphens}{url}
\usepackage{hyperref}       % 超链接
\usepackage{indentfirst}
\usepackage{geometry}
\usepackage{fancyhdr}
\usepackage{titlesec}
\usepackage{footmisc}
\usepackage{tocloft}
\usepackage{color, soul, xcolor} % 用 color 和 soul 包
  \setulcolor{red}
  \setstcolor{blue}
  \sethlcolor{yellow}
\usepackage[commandnameprefix=always,final,markup=default,authormarkup=none]{changes}
\definechangesauthor[name={JoinNico}, color=orange]{Join}

\pagestyle{plain}
\geometry{a4paper,scale=0.7}

\makeatletter
% 解决 url 超出纸面问题
\def\UrlAlphabet{%
      \do\a\do\b\do\c\do\d\do\e\do\f\do\g\do\h\do\i\do\j%
      \do\k\do\l\do\m\do\n\do\o\do\p\do\q\do\r\do\s\do\t%
      \do\u\do\v\do\w\do\x\do\y\do\z\do\A\do\B\do\C\do\D%
      \do\E\do\F\do\G\do\H\do\I\do\J\do\K\do\L\do\M\do\N%
      \do\O\do\P\do\Q\do\R\do\S\do\T\do\U\do\V\do\W\do\X%
      \do\Y\do\Z}
\def\UrlDigits{\do\1\do\2\do\3\do\4\do\5\do\6\do\7\do\8\do\9\do\0}
\g@addto@macro{\UrlBreaks}{\UrlOrds}
\g@addto@macro{\UrlBreaks}{\UrlAlphabet}
\g@addto@macro{\UrlBreaks}{\UrlDigits}
\makeatother

% % 章节标题格式：取消换页、紧凑间距、中文兼容
% \titleformat{\chapter}[block]  % 控制章节标题样式
%   {\normalfont\filcenter\bfseries\Large}  % 整体格式（居中、加粗、字号）
%   {\CTEXthechapter\quad}  % 标签格式（如“第一章”）
%   {0pt}  % 标签与标题间距
%   {\Huge}  % 标题字体大小

% % 章节间距调整（上/下间距）
% \titlespacing*{\chapter}{0pt}{10pt}{20pt}  % {左}{上}{下}{右}

% % 可选：调整 part 标题格式（如果需要）
% \titleformat{\part}[display]
%   {\normalfont\centering\bfseries\Huge}
%   {\partname\ \thepart}
%   {20pt}
%   {\Huge}
% \titlespacing*{\part}{0pt}{30pt}{40pt}


\newcommand{\BookTitle}{绝望即壤}
\newcommand{\BookSubtitle}{《魔法少女小圆》非官方批评研究}
\newcommand{\BookAuthor}{Jed A. Blue}
\newcommand{\BookTranslator}{JoinNico}
\newcommand{\BookEdition}{B6 内页样张}

\begin{document}

% 1--2: inner title and copyright page.
\frontmatter
\begin{titlepage}
  \thispagestyle{empty}
  \begin{tikzpicture}[remember picture,overlay]
    % A quiet double halo gathers the two soul gems into one crest while
    % preserving the pink/violet duality used throughout the book.
    \draw[MadokaPink,opacity=.12,line width=.35pt]
      ($(current page.north)+(-8mm,-42mm)$) circle (18mm);
    \draw[HomuraViolet,opacity=.12,line width=.35pt]
      ($(current page.north)+(8mm,-42mm)$) circle (18mm);
    \node[opacity=.98] at ($(current page.north)+(0,-42mm)$)
      {\SoulGemPair[1.82]};

    \node[
      text=BookGray,
      font=\sffamily\fontsize{6.8pt}{8pt}\selectfont
    ] at ($(current page.north)+(0,-15mm)$)
      {\addfontfeatures{LetterSpace=25}THE VERY SOIL};
    \draw[BookGray,opacity=.42,line width=.28pt]
      ($(current page.north)+(-39mm,-20mm)$)--($(current page.north)+(-8mm,-20mm)$)
      ($(current page.north)+(8mm,-20mm)$)--($(current page.north)+(39mm,-20mm)$);
    \fill[MadokaPink]
      ($(current page.north)+(-2.2mm,-20mm)$) circle (.45mm);
    \fill[HomuraViolet]
      ($(current page.north)+(2.2mm,-20mm)$) circle (.45mm);

    \node[
      text=BookBlack,
      align=center,
      font=\bookparttitlefont\fontsize{27pt}{32pt}\selectfont
    ] at ($(current page.north)+(0,-78mm)$)
      {\addfontfeatures{LetterSpace=10}\BookTitle};
    \draw[MadokaPink,line width=.65pt]
      ($(current page.north)+(-26mm,-89mm)$)--($(current page.north)+(-2mm,-89mm)$);
    \draw[HomuraViolet,line width=.65pt]
      ($(current page.north)+(2mm,-89mm)$)--($(current page.north)+(26mm,-89mm)$);
    \node[
      text=BookGray,
      align=center,
      text width=88mm,
      font=\sffamily\fontsize{9.6pt}{14pt}\selectfont
    ] at ($(current page.north)+(0,-101mm)$)
      {\BookSubtitle};
    \node[
      text=HomuraViolet,
      opacity=.72,
      font=\fontsize{7.2pt}{9pt}\selectfont
    ] at ($(current page.north)+(0,-114mm)$)
      {\runeecho{THE VERY SOIL}};

    % The credits share one baseline and read as a deliberate colophon,
    % rather than two loose lines pushed into the lower-right corner.
    \draw[BookGray,opacity=.35,line width=.28pt]
      ($(current page.south)+(-40mm,38mm)$)--($(current page.south)+(40mm,38mm)$);
    \draw[BookGray,opacity=.32,line width=.28pt]
      ($(current page.south)+(0,34mm)$)--($(current page.south)+(0,18mm)$);
    \node[
      anchor=north,
      text=MadokaPink,
      font=\sffamily\bfseries\fontsize{6.5pt}{8pt}\selectfont
    ] at ($(current page.south)+(-22mm,33mm)$)
      {\addfontfeatures{LetterSpace=10}作者 / AUTHOR};
    \node[
      anchor=north,
      text=BookBlack,
      font=\rmfamily\fontsize{10.5pt}{13pt}\selectfont
    ] at ($(current page.south)+(-22mm,26mm)$)
      {\BookAuthor};
    \node[
      anchor=north,
      text=HomuraViolet,
      font=\sffamily\bfseries\fontsize{6.5pt}{8pt}\selectfont
    ] at ($(current page.south)+(22mm,33mm)$)
      {\addfontfeatures{LetterSpace=10}译者 / TRANSLATOR};
    \node[
      anchor=north,
      text=BookBlack,
      font=\rmfamily\fontsize{10.5pt}{13pt}\selectfont
    ] at ($(current page.south)+(22mm,26mm)$)
      {\BookTranslator};
    \draw[MadokaPink,line width=.7pt]
      ($(current page.south west)+(14mm,10mm)$)--($(current page.south)+(0,10mm)$);
    \draw[HomuraViolet,line width=.7pt]
      ($(current page.south)+(0,10mm)$)--($(current page.south east)+(-14mm,10mm)$);
  \end{tikzpicture}
  \null
\end{titlepage}

\clearpage
\thispagestyle{empty}
\vspace*{\fill}
{\fontsize{8pt}{12pt}\selectfont
  \textbf{\BookTitle}\par
  \BookSubtitle\par
  \vspace{8pt}
  作者：\BookAuthor\par
  译者：\BookTranslator\par
  版本：\BookEdition\par
  印刷：\BookPrinting\par
  开本：\BookFormat\par
  % 输出：\ifbookprint 印刷版\else 屏幕阅读版\fi\par
  \vspace{8pt}
  译文与排版 © 2026 JoinNico。保留所有权利；未经许可，不得复制、转载、改编或用于商业用途。\par
  字体与外部素材遵循各自许可文件。原著及相关作品的权利归各自权利人所有。\par
  本书为非官方翻译评论作品，与原作者、出版方及相关权利方无隶属或授权关系。\par
  \vspace{8pt}
  排版系统：\LaTeX\par
}
\clearpage


% 3: contents specimen.
\thispagestyle{empty}
\begin{center}
  {\sffamily\bfseries\fontsize{18pt}{22pt}\selectfont 目录\par}
\end{center}
\vspace{12pt}
{\sffamily\bfseries\large\color{MadokaPink}\hfill 第一部分\quad 电视动画\hfill\par}
\vspace{9pt}
\noindent 第一章\quad 一股秘溪潺潺流淌\dotfill 3\par
\noindent 第二章\quad 并非一时兴起\dotfill 9\par
\noindent 第九章\quad 逐渐而间接地获得政治意义\dotfill 67\par
\vspace{12pt}
{\sffamily\bfseries\large\color{HomuraViolet}\hfill 第三部分\quad 叛逆电影\hfill\par}
\vspace{9pt}
\noindent 第二十章\quad 叛逆的焰\dotfill 151\par
\clearpage

% 4: Part page specimen.
\thispagestyle{empty}
\begin{tikzpicture}[remember picture,overlay]
  \node at ($(current page.center)+(\BookRectoTextCenterShift,17mm)$) {\HopeMotif[2.5]};
\end{tikzpicture}\par
\vspace*{69mm}
% This diagnostic page is physically even; shift its text to simulate the
% recto text block used by every real Part page.
\noindent\hspace*{7mm}\makebox[0pt][l]{%
  \begin{minipage}[t]{\textwidth}
    \begin{center}
      {\sffamily\bfseries\fontsize{9pt}{11pt}\selectfont\color{BookGray}PART I\par}
      \vspace{5pt}
      {\sffamily\bfseries\fontsize{22pt}{28pt}\selectfont 电视动画\par}
      \vspace{8pt}
      {\color{HomuraViolet}\runeecho{PART I}\par}
    \end{center}
  \end{minipage}%
}\par
\clearpage

% 5: diagnostic plate page only; publication builds do not include it.
\thispagestyle{empty}
\begin{tikzpicture}[remember picture,overlay]
  \fill[BookPale] ($(current page.south west)+(8mm,8mm)$)
    rectangle ($(current page.north east)+(-8mm,-8mm)$);
  \draw[MadokaPink,dashed,line width=.5pt]
    ($(current page.south west)+(11mm,11mm)$)
    rectangle ($(current page.north east)+(-11mm,-11mm)$);
  \node[align=center,text=BookGray,font=\sffamily] at (current page.center)
    {铜版彩插诊断页\par\small 外部画稿到齐后启用};
\end{tikzpicture}
\clearpage

% 6: ordinary chapter opener.
\pagenumbering{arabic}
\setcounter{chapter}{0}
\chapter{一股秘溪潺潺流淌}
\chaptersubtitle{《似乎在梦里见过，那样……》}
《魔法少女小圆》的剧情可划分为三个相互交织又各有侧重的篇章。每个篇章聚焦于特定角色，并深入探讨一个核心主题。这里用于确认章首标题、集名副标题与正文首段之间的节奏。
\clearpage

% 7: long chapter-opening note.
\thispagestyle{chapteropening}
\begin{chapteropeningnote}
\textbf{原文}\quad Perhaps hopelessness is the very soil that nourishes human hope; perhaps one could never find sense in life without first experiencing its absurdity.

\textbf{译文}\quad 也许绝望正是滋养人类希望的土壤；也许一个人只有先经历生命的荒谬，才能从中发现意义。

本框用于章标题出处、译文与必要说明。它可以跨页，但不再挤占普通页脚，也不会和文献引用编号混淆。
\end{chapteropeningnote}
章首题记之后，正文恢复八磅字体与十三磅行距。粉色边线提供视觉识别，文字信息本身仍由标签与位置表达，因此灰度输出不会丢失层级。
\clearpage

% 8: dense body, short footnotes, Latin and citation specimens.
\pagestyle{bookbody}
\markboth{第一部分\quad 电视动画}{第一章\quad 一股秘溪潺潺流淌}
\noindent
动画伊始，以定格动画呈现、非传统动画形式的纸幕缓缓升起。艺术与非艺术之别，关键在于“框架”\translatornote{这里是一条较长的译者说明，用于测试页脚字号、悬挂缩进、连续两行的基线，以及它和下一条脚注之间的垂直距离是否一致。}。一篇没有框架的故事只是谎言；一幅没有框架的画作只是涂鸦，正是框架定义了艺术。

The Latin companion face is Libertinus Serif. It must remain readable beside Chinese prose, preserve quotation marks and italics, and avoid the broken spacing visible in the previous A4 edition.

这种诠释颇为契合，但让我们先回到剧情。开篇的视觉母题显得平静，随后纸幕、裂纹与不稳定的空间逐步侵入日常世界\footnote{普通作者注采用七磅字和九点五磅行距。}。
\clearpage

% 9: interlude specimen.
\renewcommand{\BookSpecialKicker}{INTERLUDE I}
\chapter*{间幕 一\quad 完美一刻}
\renewcommand{\BookSpecialKicker}{}
在动画第二集中，墙上出现的德语涂鸦直接引用了《浮士德》的文本。间幕沿用章首骨架，但使用独立英文标签，使其在目录与正文节奏中都能被立即识别。
\clearpage

% 10: bibliography and URL specimen.
\thispagestyle{bookbody}
\chapter*{参考文献}
\nocite{aristotle2009poetics,magicaquartet2011production-note}
\printbibliography[heading=none]

\end{document}
